% !TEX program = pdflatex
\documentclass[11pt]{article}
\usepackage{support-notes}
\title{Failure of constructibility for bounded Ext at fixed graded cohomology}
\author{\workingauthor}\date{\coursedate}
\begin{document}\maketitle
\begin{abstract}
In a flat six-dimensional two-term dg family on the complex multiplicative
group, a fixed perfect cone and a uniformly presented detector have
bounded total Ext exactly away from the roots of unity. Equality of two
named detector profiles has the same nonconstructible locus. Even the
abstract poset distinguishes the two sets: it has an atom precisely at
infinite-order parameters. This occurs
with constant graded cohomology and nonzero cubic multiplication
throughout. We also compute the whole profile order of the infinite-order
quantum plane on the Koszul-dual side. These statements are proved.
\end{abstract}

\section{A fixed pair in a flat family}
Let $\Lambda=k[q,q^{-1}]$. Over $\Lambda$ take $D^0$ with basis
$1,a,b,h$, $D^{-1}$ with basis $e,w$, and define
\[
 a^2=h,\qquad e b=w,\qquad b e=q w,\qquad d e=h .
\]
All other radical products and differentials vanish.
The table is associative and Leibniz because triple products vanish
and $h$ kills the radical. The complex splits as free cohomology
plus a contractible pair, and every fibre has the same graded
cohomology algebra with square-zero radical. The two values of the
ternary operation are $q w$ and $-w$. At least one is nonzero and
normalized coboundaries on reduced inputs vanish.

The residue-field dual in every field-valued fibre is
$A_q=k\langle x,z\rangle/(x^2z-qzx^2)$ in generator degree one.
This calculation works over $\Lambda$ as well. Before transfer the
strict reduced bar dual is free on $x,z,u$ in degree one and $v,y$ in
degree two, with signs chosen so that
$du=v-x^2$, $dy=vz-qzv$ and $dx=dz=dv=0$.
In the suspended bar convention the terms on $[e|b]$ and $[b|e]$
have opposite multiplication signs. Replacing $v$ by $v-x^2$ and $y$
by $y-(uz+qzu)$ leaves $dy=x^2z-qzx^2$ on the remaining variables;
the extra pair contracts. Its deletion is a quasi-isomorphism by the
alternating word decomposition into tensor products with contractible
free pairs. Filtering by the expansion $y\mapsto xxz$, the independent
disjoint occurrences of $xxz$ give contractible two-choice complexes
with unit leading coefficient. Thus projection onto the displayed
quotient is a quasi-isomorphism. The normal words avoiding $xxz$ form
a degreewise free $\Lambda$-basis. All these maps commute with base
change. These are also the relative endomorphisms of the augmentation:
the augmented relative bar resolution has terms free over $D$, since
the radical is free over $\Lambda$; it computes derived Hom out of
$\Lambda$ (its usual length filtration is semi-free after resolving the
split underlying chain pairs). Equivalently one may use the explicit
bar endomorphism action by deconcatenation and augmentation. The finite
cones below consequently define uniform objects of $\thick_D(\Lambda)$
on the dg side. Only the fibrewise equivalence to $D_{\fd}$ is used.

For a scalar $\lambda$ define the graded right module $M_\lambda$,
free of rank one over the base in each nonnegative degree, by
\[
 e_i x=e_{i+1},\quad e_{2j}z=\lambda q^j e_{2j+1},\quad
 e_{2j+1}z=0 .
\]
The relation follows from multiplication by $q$ every two steps. There
is a free presentation $0\to A[-1]\oplus A[-2]\to A\to M_\lambda\to0$
with first arrow $(z-\lambda x,\,xz)$. Its leading normal prefixes $z,xz$
are disjoint and the lower term has one fewer $z$, proving independence.
Every normal word starting in $z$ or $xz$ reduces in the cokernel by
these relations and induction on the number of $z$'s to a multiple of
a power of $x$. The map to $M_\lambda$ proves these powers are free
over the base. This proof gives exactness over $\Lambda$ and in all fibres.
Use the fixed cone $C=\Cone(z-x:A[-1]\to A)$.

\begin{thm}\label{jump}
Over $k=\mathbb C$, $C\in\UU_{M_1}$ exactly at parameters $q$
of infinite multiplicative order. Furthermore
$\UU_{M_1}=\UU_{M_2}$ exactly at the same parameters. Neither locus
is Zariski constructible, open or closed.
\end{thm}
\begin{proof}
In high degrees the map computing Hom from $C$ to $M_1$ is the scalar
$q^j-1$ in even degrees $2j$ and $-1$ in odd degrees. The kernel or
cokernel contributes to the high groups of the cone whenever the scalar
is zero. The even scalar vanishes infinitely often exactly at roots of
unity. For $M_2$ at any root of unity the scalar $2q^j-1$ never vanishes.
Thus $C$ separates the two profiles at every such parameter.

For infinite-order $q$, localize at the normal element $x^2$. Both
targets are eventually invertible there. As in Lemmas 2 and 3 of
\othernote{normal_localisations}{Bounded profiles and finite-dimensional
points of universal localisations}, any perfect source on the localization
has a finite degree-zero matrix presentation over the free algebra
on $U=zx^{-1}, V=xzx^{-2}$, and its high Hom into either target is
computed there. This step uses no periodicity assumption. On the
degree-$2j$ and degree-$(2j+1)$ components the coefficients specialize
at $(\lambda q^j,0)$ and $(0,\lambda q^{j+1})$. An unequal rectangular
matrix cannot pass. For a square matrix $H$, the two required tests
are that $\det H(t,0)$ and $\det H(0,t)$ are nonzero polynomials in
one commuting indeterminate. Indeed a nonzero polynomial has only finitely
many roots and an infinite-order progression with $\lambda\ne0$ has
distinct terms. A zero polynomial fails at every corresponding degree.
The criteria coincide at $\lambda=1,2$. Finally roots of unity form an
infinite subset with infinite complement in the complex torus; a
constructible subset of an irreducible algebraic curve is finite or
cofinite. This proves all assertions.
\end{proof}

For the membership assertion alone the field can be arbitrary and one
asks whether its chosen unit is torsion. The geometric nonconstructibility
assertion is on complex points. No single exponent bound detects eventual
vanishing in this family. Each individual degree remains a finite rank
test. This result refutes general constructibility of fixed-pair boundedness
and of equality for uniformly given detector profiles. A stronger, intrinsic statement follows.

\begin{thm}\label{atom}
For all $n,r\ge1$ with $g=nr\ge2$, the profile poset for
$A_{n,r,qI}$ over $\mathbb C$ has an atom (a minimal element strictly above
its bottom) if and only if $q$ has infinite order. In particular this
isomorphism-invariant property has a nonconstructible locus in every
constant-cohomology two-term family with $n=2$, including the family above.
\end{thm}
\begin{proof}
The bottom is $\UU_A$: a perfect source $P$ is in it iff the left
perfect dual $Y=P^\vee$ has finite cohomology. This suffices for
bounded tensor with every perfect detector. If $q$ has finite order,
use the full two-row theorem in free rank $g$. Each finite first-row label
can be enlarged by a new simple orbit, strictly decreasing its profile
while retaining the source $G=A/(x^n)$, which is outside the bottom.
For a nonempty second-row domain $O(H)$ choose a simple tuple
$(S_1,\ldots,S_g)$ in it, of dimension $d$. The tuples
$(S_1+cI,S_2,\ldots,S_g)$ remain simple for
every $c\in\mathbb C$. Every one of the finitely many twisted
invertibility conditions has a determinant polynomial in $c$ nonzero
at zero, so all hold on a cofinite set. The trace of $S_1+cI$ distinguishes
infinitely many isomorphism classes, since $d\ne0$ in $\mathbb C$.
Their finite twist orbits are still infinite in number. Pick two
distinct such orbits. By the Chinese remainder argument in the
localisation theorem a polynomial $f$ kills one orbit and is identity
on the other. The domain of $H\oplus f$ is a proper nonempty subdomain.
Consequently no profile covers the bottom in these fibres.

Now let $q$ have infinite order. Take the tame detector whose localized
simple is the one-dimensional zero tuple; write $M_0$ for its perfect lift.
Its profile contains $G$ and omits $A$, so it is proper and above
the bottom (the free resolution of $G$ has dual with infinite cokernel
$G$ up to shift). We show that it lies inside every eventually invertible
detector profile. A source in $\UU_{M_0}$ has a localized presentation
matrix $H$ that evaluates bijectively at the zero tuple, so it is
square and $H(0)$ is invertible. For any finite tuple on the
localization the $n$th power of the twist scales all $g$ operators
by $q$. The twists thus lie in $n$ progressions of scalar multiples
of fixed tuples (using $q^{-1}$ for the reverse convention). In each progression
the evaluation determinant is a polynomial in $q^j$ with nonzero
constant, since that constant is a power of $\det H(0)$.
It has finitely many exceptional degrees. This proves the inclusion
for every such detector, including those with zero localization.

For a detector failing eventual invertibility, the relative cut bound
forces $Y=P^\vee$ to be eventually invertible for each source in its
profile. If this source also lies in $\UU_{M_0}$, tensor the localized
free resolution of the one-dimensional zero tuple (ranks $g$ and one)
with $Y_h$. Whenever $H^0(Y_h)\ne0$, the two evaluated terms have
different dimensions in every degree, so the localized tensor cannot
be bounded above. Thus $H^0(Y_h)=0$. Eventual invertibility and the
degree-one unit on the chart give finite total cohomology of $Y$.
The intersection with $\UU_{M_0}$ is therefore exactly the bottom.
No profile lies strictly between the bottom and $\UU_{M_0}$.
For $n=2$ use the strict table $b_i e=q w_i$, $e b_i=w_i$,
$a^2=h$, $de=h$. The bar computation above works independently for
each $b_i$ and the underlying cohomology stays free and constant.
\end{proof}

\section{A boundary calculation in free rank one}
Let $A=k\langle x,z\rangle/(xz-qzx)$ and let $q\in k^\times$ have
infinite order. This is the period-one case of the normal family.
Here the open ring is $k[U]$, $U=zx^{-1}$, with $\sigma U=qU$.
In the finite-order case its classification gives the orbit version of
finite subsets of the projective line plus a bottom. Infinite order
has a different order already in this small rank.

\begin{thm}\label{six}
For this $A$ and any infinite-order $q$ over any field, there are
exactly six bounded profiles. From the top down the covering relations
are
\[
 \mathrm{top}\ >\ T_{\rm mov}\ >
 \{T_0,B_{\rm all}\}\ >\ B_{\rm mov}\ >\ \mathrm{bottom},
\]
where the two middle entries are incomparable and every adjacent-level
pair in the display is a cover.
\end{thm}
\begin{proof}
The normal-element test applies with fibre $k[z]$. Eventually invertible
detectors have finite degree-zero localized modules over $k[U]$. All
nonzero closed points have infinite scaling orbits: a nonzero algebraic
root of a monic irreducible with finite orbit would be permuted by a
power of $q$, forcing a further power to equal one. This does not use
separability. For a localized square presentation $H$, put $f=\det H$.
On a finite primary module the map is invertible precisely when $f$
is nonzero at the corresponding point. Along the twists of a nonzero
point, a fixed nonzero polynomial can fail only finitely many times,
in either direction of twisting. At the zero point it must pass every
time. Thus a zero target on the localization gives the full profile;
a nonzero target supported away from zero gives $T_{\rm mov}$, requiring
a square presentation with $f\ne0$; a target with a zero primary factor
gives $T_0$, requiring also $f(0)\ne0$. These conditions concern one
injective presentation from the chart lemma; the underlying acyclicity
test makes them independent of its choice. The source $A$ separates
the full profile, and the cone of $z$ distinguishes the other two.

For a detector failing eventual invertibility, the source dual must
be eventually invertible. Its finite degree-zero localized module $N$
is then tested against the presentation of the detector. A rectangular
unequal-rank or identically singular presentation requires $N=0$
(the bottom). A nonzero determinant vanishing at zero allows just those
$N$ supported away from zero, including $0$ ($B_{\rm mov}$).
Nonvanishing at zero allows all $N$ ($B_{\rm all}$). $A$, $G\oplus
\Cone(z)$ and $G=A/(x)$ realise these choices. Duals of the left
one-dimensional modules at 1 and 0 separate the three.

The source $G$ is in each first-row profile and no second-row profile.
For a possible reverse inclusion, take finite right and left modules
on the localization. Tensoring their primary resolutions in each twist
can have nonzero high cohomology infinitely often precisely when
both supports contain zero. Two finite sets of nonzero closed points
can coincide under only finitely many twists because all their orbits
are infinite. Hence $B_{\rm all}\subset T_{\rm mov}$ and
$B_{\rm mov}\subset T_0$, whereas $B_{\rm all}\not\subseteq T_0$.
These are exactly the stated comparisons. The closed and open tests are
valid even though no positive power of $x$ is central here.
\end{proof}

\paragraph{Prior work and scope.}
The orbit mechanism has a close predecessor in the four-dimensional
quantum exterior algebra: the modules studied by
\href{https://www.cambridge.org/core/services/aop-cambridge-core/content/view/26516CA33D32E5B50A8C291AEE0EB145/S1446788716000331a.pdf/ext-finite-modules-for-weakly-symmetric-algebras-with-radical-cube-zero.pdf}{Erdmann}
have syzygies which scale the parameter; the paper also cites Schulz
and earlier explicit computations. Bergh studies vanishing symmetry
for \href{https://arxiv.org/abs/0811.4309}{quantum complete intersections}.
Theorem~\ref{six} organizes this familiar rank-one phenomenon into all
bounded profiles of the derived category in the stated convention.
Theorem~\ref{jump} transports it to a family whose graded cohomology
algebra is fixed, and proves equality of profiles by testing every
perfect source. The orbit phenomenon itself has a substantial history.
For arbitrary twists in characteristic zero, the companion note
\othernote{orbit_closures}{Orbit closures and bounded profiles}
gives a geometric description using generic matrix ranks on every
irreducible component. Intrinsic equality of its upper labels without
quantifying over matrix tests remains open in free rank at least two.
\end{document}
